\input{../tex-inputs/latex-leseansicht-vorspann}

\section[Emil Marriot an Arthur Schnitzler, 9. 2. 1903]{L04338 Emil Marriot an Arthur Schnitzler, 9. 2. 1903}
\nopagebreak\mylabel{L04338v}
\rehead{ }\normalsize\beginnumbering\briefempfaengerindex{Schnitzler, Arthur@\textsc{Schnitzler, Arthur}!zzzMarriot, Emil@\emph{von Emil Marriot}!1903-02-091@{9. 2. 1903}|(be}
\toendnotes[C]{\smallbreak\pagebreak[2]}
\correspDesc{Versand  durch Emil Marriot am 9. 2. 1903 in Wien
\newline{}Erhalt  durch Arthur Schnitzler im Zeitraum [10. 2. 1903 – 14. 2. 1903?] in Wien}\toendnotes[C]{\smallbreak}
\Standort{DLA, A:Schnitzler, HS.1985.1.4008.}
\physDesc{Brief, 1 Blatt, 2 Seiten, 1785 Zeichen
\newline{}Handschrift: schwarze Tinte, lateinische Kurrent
\newline{}Schnitzler: 1) mit Bleistift Vermerk: »\textsc{Marriot}«  2) mit rotem Buntstift eine Unterstreichung}\toendnotes[C]{\smallbreak}
\pstart
           \raggedleft{}{\pb}Wien\oindex{Wien@\textbf{Wien}, \emph{Verwaltungsgebiet}|pw},
               9. II. 1903.\pend
           
\pstart
           \raggedleft{}II/2, Schüttelstrasse 31\oindex{Wien@\textbf{Wien}!II., Leopoldstadt@\textbf{II., Leopoldstadt}!Schüttelstraße 31@\textbf{Schüttelstraße 31}, \emph{Straße}|pw}.\pend
           
\pstart{}Sehr geehrter Herr Doktor!\pend\vspace{0.5em}
\pstart
           Vielleicht – wahrscheinlich – wissen Sie bereits,
               dass Herr \label{K_L04338-1v}\edtext{Harden\pwindex{Harden, Maximilian 20.\,10.\,1861 Berlin – 30.\,10.\,1927 Montana@\textsc{Harden, Maximilian} (20.\,10.\,1861 Berlin – 30.\,10.\,1927 Montana), \emph{Schriftsteller, Publizist}|pw} am 27. u 28. d. Vorträge}{\lemma{\textnormal{\emph{Harden … Vorträge}}}\Cendnote{\textnormal{Die hier
                  geplanten Vorträge wurden von Harden\pwindex{Harden, Maximilian 20.\,10.\,1861 Berlin – 30.\,10.\,1927 Montana@\textsc{Harden, Maximilian} (20.\,10.\,1861 Berlin – 30.\,10.\,1927 Montana), \emph{Schriftsteller, Publizist}|pwk} abgesagt und fanden nicht statt.}}}\label{K_L04338-1}
      hier halten wird. Für den ersten, am 27.,
               der zum Besten des socialen Hilfsvereins\orgindex{Frauenverein für sociale Hilfsthätigkeit@Frauenverein für sociale Hilfsthätigkeit|pw} im Saal der Ingenieure u Architekten\oindex{Wien@\textbf{Wien}!I., Innere Stadt@\textbf{I., Innere Stadt}!Österreichischer Ingenieur- und Architektenverein@\textbf{Österreichischer Ingenieur- und Architektenverein}|pw} stattfinden wird, bin
      ich gewissemassen mit verantwortlich, und
      deshalb wünsche ich sehr, er möchte befriedigend
      ausfallen. Darunter verstehe ich hauptsächlich
      ein intelligentes Publikum und ich möchte bei
      Zeiten dafür sorgen, dass der Saal\oindex{Wien@\textbf{Wien}!I., Innere Stadt@\textbf{I., Innere Stadt}!Österreichischer Ingenieur- und Architektenverein@\textbf{Österreichischer Ingenieur- und Architektenverein}|pwv} vorwiegend
      mit solchem gefüllt werde. Vielleicht wollen
      sie die Güte haben, mich dabei zu unterstützen?
      Die Karten werden in der Manz’schen Buchhandlung\orgindex{Manz’sche Verlags- und Universitätsbuchhandlung@Manz’sche Verlags- und Universitätsbuchhandlung|pw} zu haben sein. Doch nehmen die
      Damen vom Comité jetzt schon Vormerkungen
      entgegen oder senden auf Wunsch die Billets;
      und es empfiehlt sich, sich bei Zeiten zu
      versorgen, da wir natürlich Denen, die zuerstkommen, die besseren Plätze geben können. Wenn
      Sie also für sich oder Angehörige und
      Freunde Bedarf haben, bitte ich, sich an
      Fräulein Helene Migerka\pwindex{Migerka, Helene 13.\,9.\,1867 Brünn – 26.\,3.\,1928 Graz@\textsc{Migerka, Helene} (13.\,9.\,1867 Brünn – 26.\,3.\,1928 Graz), \emph{Schriftstellerin}|pw}, Wien, II/2, Czerningasse 7\oindex{Wien@\textbf{Wien}!II., Leopoldstadt@\textbf{II., Leopoldstadt}!Czerningasse 7@\textbf{Czerningasse 7}, \emph{Wohngebäude}|pw} zu wenden. Der Preis der
               Karten ist 10, 8, 6, 4, 3 und 2 Kronen \introOben{}(Stehplätze.)\introOben{}\pend
           
\pstart
           {\pb}Nach dem Vortrag wollen wir uns
               mit Herrn Harden\pwindex{Harden, Maximilian 20.\,10.\,1861 Berlin – 30.\,10.\,1927 Montana@\textsc{Harden, Maximilian} (20.\,10.\,1861 Berlin – 30.\,10.\,1927 Montana), \emph{Schriftsteller, Publizist}|pw} in einem Hôtel zusa{\geminationm}enfinden und es würde ihn gewis besonders
      freuen, wenn auch Sie dabei wären. Ich
      würde mich ebenfalls freuen, Sie bei dieser
      Gelegenheit wiederzusehen. Das Nähere: Hôtel,
      Stunde, Couvert u. s. w. ist noch nicht vereinbart. Sobald wir ungefähr wissen, auf wie
      viele Theilnehmer wir rechnen können, werden
      wir uns über alle Details einigen.\pend
           
\pstart
           Der »\label{K_L04338-2v}\edtext{Schleier du Beatrice\pwindex{Schnitzler, Arthur 15. 5. 1862 Wien – 21. 10. 1931 ebd.@\textsc{Schnitzler, Arthur} (15. 5. 1862 Wien – 21. 10. 1931 ebd.), \emph{Schriftsteller, Mediziner}!Schleier der Beatrice. Schauspiel in fünf Akten@\strich\emph{Der Schleier der Beatrice. Schauspiel in fünf Akten}|pw}« kommt also
               doch heraus? Ich freue mich darüber. Vielleicht fällt die Première\eventindex{Deutsches Theater Berlin@\textbf{Deutsches Theater Berlin}!Premiere von Der Schleier der Beatrice, 7.3.1903@Premiere von Der Schleier der Beatrice, 7.3.1903|pwv}}{\lemma{\textnormal{\emph{Schleier … Première}}}\Cendnote{\textnormal{Die Premiere von \emph{Der Schleier der Beatrice. Schauspiel in fünf Akten}\pwindex{Schnitzler, Arthur 15. 5. 1862 Wien – 21. 10. 1931 ebd.@\textsc{Schnitzler, Arthur} (15. 5. 1862 Wien – 21. 10. 1931 ebd.), \emph{Schriftsteller, Mediziner}!Schleier der Beatrice. Schauspiel in fünf Akten@\strich\emph{Der Schleier der Beatrice. Schauspiel in fünf Akten}|pwk} von Arthur Schnitzler\eventindex{Deutsches Theater Berlin@\textbf{Deutsches Theater Berlin}!Premiere von Der Schleier der Beatrice, 7.3.1903@Premiere von Der Schleier der Beatrice, 7.3.1903|pwk} fand am 7. 3. 1903 am \emph{Deutsches Theater Berlin}\orgindex{Deutsches Theater Berlin@Deutsches Theater Berlin|pwk} statt.}}}\label{K_L04338-2} übrigens gerade mit
               Hardens\pwindex{Harden, Maximilian 20.\,10.\,1861 Berlin – 30.\,10.\,1927 Montana@\textsc{Harden, Maximilian} (20.\,10.\,1861 Berlin – 30.\,10.\,1927 Montana), \emph{Schriftsteller, Publizist}|pw} kommen zusa{\geminationm}en. In diesem Fall
      müssten wir freilich auf Ihre Theilnahme verzichten.\pend
           
\pstart
           Mit verbindlichen Grüssen{\\[\baselineskip]}Ihre ganz ergebene{\\[\baselineskip]}\spacefill\mbox{E. Marriot.}\pend
           \leftskip=0em{}\selectlanguage{ngerman}\endnumbering\briefempfaengerindex{Schnitzler, Arthur@\textsc{Schnitzler, Arthur}!zzzMarriot, Emil@\emph{von Emil Marriot}!1903-02-091@{9. 2. 1903}|)be}\mylabel{L04338h}
\begin{anhang}
\end{anhang}\newcommand{\dateiname}{L04338}\newcommand{\titel}{Emil Marriot an Arthur Schnitzler, 9. 2. 1903}\newcommand{\editorInnen}{Herausgegeben von Jahnke, SelmaMüller, Martin Anton}\input{../tex-inputs/latex-leseansicht-abspann}
